\section{Agent Orchestration: סוכן יחיד מול ריבוי סוכנים}

כמשימות מתמערות, נדרשת החלטה ארכיטקטונית: סוכן אחד עם מגוון \eng{Skills}
(``כובעים'') או מספר סוכנים במקביל, כל אחד עם התמחות. ההרצאה דנה
ביתרונות, חסרונות והצורך ב-\textbf{אורכוסטרטור}.

\subsection{סוכן יחיד, ריבוי כובעים}

סוכן אחד בטרמינל יכול לטעון \eng{Skill} שונה לכל משימה --- היום מומחה
לטבלאות, מחר מומחה לתמונות. היתרונות:
\begin{itemize}[rightmargin=2em]
  \item פשטות ניהולית --- מעקב אחרי תהליך אחד.
  \item עומס חישובי נמוך --- רק instance אחד פעיל.
  \item קל לדבג --- יודעים היכן קרתה שגיאה.
\end{itemize}

החיסרון: ביצוע \textbf{סדרתי}. אם בדיקת טבלאות לוקחת דקה ובדיקת תמונות
דקה נוספת --- סה''כ שתי דקות.

\subsection{ריבוי סוכנים, כובע אחד לכל אחד}

פתיחת מספר טרמינלים (או תהליכים) --- כל סוכן עם \eng{Skill} ייעודי.
יתרון מרכזי: \textbf{מקביליות}. טבלאות בעמוד 2 ותמונות בעמוד 1 יכולים
להיבדק בו-זמנית על אותו \texttt{main.tex}, כל עוד הם לא משנים את אותן
שורות.

דוגמה מההרצאה: 50 משימות של דקה כל אחת ---
\begin{itemize}[rightmargin=2em]
  \item סוכן יחיד: 50 דקות.
  \item 50 סוכנים במקביל: כדקה אחת (בתיאוריה).
\end{itemize}

\subsection{חסרונות ריבוי סוכנים}

\begin{itemize}[rightmargin=2em]
  \item מורכבות ניהולית --- קשה לעקוב אחרי 50 תהליכים.
  \item סיכון לולאות וצריכת טוקנים לא מבוקרת.
  \item מגבלת חומרה --- המחשב ``נחנק'' מעומס.
  \item התנגשויות כתיבה --- נעול קבצים ברמת מערכת ההפעלה.
\end{itemize}

\subsection{תפקיד האורכוסטרטור}

כשמפעילים מספר סוכנים, \textbf{חובה} סוכן-על (\eng{Orchestrator}):
\begin{enumerate}[rightmargin=2em]
  \item מחלק משימות לסוכני משנה.
  \item אוסף תוצאות ומרכיב מסמך שלם.
  \item מפקח על תקציב טוקנים וזמן.
  \item מטפל בהתנגשויות (תור כתיבה לקובץ).
\end{enumerate}

\begin{notebox}
במטלה 02 (תרגום רב-שלבי) יושמה אורכוסטרציה: שלושה סוכנים ברצף עם
\texttt{run\_pipeline.py} כמנהל. אותו עיקרון --- עם אפשרות למקביליות
בשלב ה-\eng{QA}.
\end{notebox}

\agenttable{
\caption{השוואת מודלי אורכוסטרציה}
\label{tab:orchestration}
\begin{tabular}{@{}p{3.5cm}p{4cm}p{4cm}@{}}
\toprule
\textbf{מאפיין} & \textbf{סוכן יחיד} & \textbf{ריבוי סוכנים} \\
\midrule
מהירות & איטית & מהירה (מקביל) \\
ניהול & פשוט & מורכב \\
משאבים & נמוכים & גבוהים \\
דיבוג & קל & קשה \\
צורך באורכוסטרטור & אופציונלי & חובה \\
\bottomrule
\end{tabular}
}

\subsection{המלצה מעשית}

המרצה ממליץ על \textbf{תמהיל}: לא 50 סוכנים בבת אחת, ולא סוכן יחיד
לכל דבר. לדוגמה במטלה 04:
\begin{itemize}[rightmargin=2em]
  \item סוכן תוכן --- כותב פרקים לפי תמלול.
  \item סוכן עיצוב --- מתחזק \eng{CLS}.
  \item סוכן \eng{QA} --- מריץ קומפילציה ובודק \eng{PDF}.
\end{itemize}

\begin{insightbox}
שאלה קלאסית למבחן (לפי ההרצאה): זמן ביצוע 50 משימות של דקה ---
סוכן יחיד לעומת 50 סוכנים. הבנת המקביליות היא מפתח להטמעה ארגונית.
\end{insightbox}
